\section{Test Environment}

The pipeline was exercised on the Docker Compose stack described in Section~\ref{ch:proposed-solution} (Kafka 2-node cluster, Spark Standalone with 1 master and 2 workers, MinIO for local object storage), run on a single local machine.

% TODO: insert host machine hardware specs (CPU, RAM, disk) once measured/confirmed.

\section{Test Cases}

Two categories of tests were used to validate the pipeline:

\begin{itemize}
    \item \textbf{Correctness, on a small dataset.} A small batch of Eventsim events is run through the full path -- Kafka, Spark Structured Streaming, Hudi bronze, dbt star schema and wide table -- and the resulting rows in \texttt{fact\_streams} / \texttt{wide\_streams} are checked against the source events for row counts, key fields, and referential integrity against the dimension tables.
    \item \textbf{Runtime and throughput, on a larger volume.} Eventsim is left running for an extended period to build up a larger event volume, and the streaming job's micro-batch processing time and Kafka consumer lag are observed under sustained load to check that the pipeline keeps up with the producer on the available 4-core cluster.
\end{itemize}

\section{Monitoring as Qualitative Evidence}

The Grafana dashboards backed by Prometheus provide qualitative evidence of the pipeline's throughput and health under load, rather than a formal benchmark.

\begin{figure}[H]
    \centering
    \includegraphics[width=0.7\textwidth]{grafana_kafka.png}
    \caption{Grafana Kafka Lag dashboard: consumer-group lag and consumption rate per topic.}
    \label{fig:grafana-kafka}
\end{figure}

\begin{figure}[H]
    \centering
    \includegraphics[width=0.7\textwidth]{grafana_spark.png}
    \caption{Grafana Spark dashboard: streaming/batch job metrics, JVM heap, and per-container CPU/memory.}
    \label{fig:grafana-spark}
\end{figure}

\section{Scalability}

Table~\ref{tab:scalability} is a placeholder for measured scalability results (dataset size vs. cluster size vs. runtime). No such measurements exist in this repository at the time of writing (\texttt{results/} is not present), so no numbers are reported here to avoid fabricating results.

\begin{table}[H]
    \centering
    \caption{Scalability results (to be measured).}
    \label{tab:scalability}
    \begin{tabular}{|l|l|l|l|}
        \hline
        \textbf{Dataset} & \textbf{Nodes} & \textbf{Edges} & \textbf{Runtime} \\
        \hline
        % TODO: insert measured results
         &  &  & \\
        \hline
    \end{tabular}
\end{table}
